%original notes, unpublished, by wdjoyner
\documentclass[12pt]{article}
\usepackage{amssymb}
\usepackage{amsthm}
\usepackage{amsgen}
\usepackage{amsfonts}
\usepackage{amsmath}

\newtheorem{theorem}{Theorem} 
\newtheorem{corollary}[theorem]{Corollary} 
\newtheorem{conjecture}[theorem]{Conjecture} 
\newtheorem{lemma}[theorem]{Lemma} 
\newtheorem{proposition}[theorem]{Proposition} 
\newtheorem{definition}[theorem]{Definition} 
\newtheorem{example}[theorem]{Example} 
\newtheorem{exercise}[theorem]{Exercise} 
\newtheorem{axiom}{Axiom} 
\newtheorem{remark}{Remark} 
\def\fff{\mathbb{F}}
\def\aaa{\mathbb{A}}
\def\qqq{\mathbb{Q}}
\def\rrr{\mathbb{R}}
\def\ccc{\mathbb{C}}
\def\zzz{\mathbb{Z}}
\def\hhh{{\cal H}}
\def\pf{{\bf proof}:\ }
\def\qed{$\Box$}


\begin{document}


\author{David Joyner}
\title{Notes on trace formulas for finite groups\thanks{Notes 
written summer 1996 visiting Stanford Univ.. 
Some details added June, November 1997, August 2000, March 2003}
}
\date{3-3-2003}
\maketitle

\tableofcontents

\vskip .4in

We will write down some analogs of the ``global" (or Selberg)
trace formula and the ``global" trace formula of Arthur
in the case of finite groups. We shall only assume a basic
familiarity with groups and representations. For the sake
of drawing an analogy, we shall often write sums as
integrals.

\textbf{Notation}: Let $G$ be a finite group, $G_{*}$ the set of conjugacy
classes, $G^{*}$ the set of equivalence classes of irreducible complex
representations. We denote a conjugacy class in $G$ by
\[
\{\gamma \}=\{g^{-1}\gamma g\;|\;g\in G\},\;\;\;\;\;\gamma \in G.
\]
We often identify $G_{*}$ (resp., $G^{*}$ ) with a complete set of
representatives of these classes. Let $C(X)$ denote the vector space of
complex-valued functions on a finite set $X$. Let $dx$ denote the counting
measure on $G$.

The vector space $C(G_{*})$ is called the space of \textbf{class functions}
on $G$.

\section{Orbital integrals and characters}


We can construct a class function
from any function $f\in C(G)$ by summing over
all the conjugate elements:
Define the \textbf{orbital integral} by 
\begin{eqnarray*}
O &:&C(G)\rightarrow C(G_{*}), \\
O(f,\gamma ) &=&\int_Gf(x^{-1}\gamma x)\, dx,\ \ \ \gamma\in G.
\end{eqnarray*}
Here $dg$ denotes the counting measure.
For $\pi \in G^{*}$, let $\mathrm{tr}$\medskip\ $\pi \in C(G_{*})$ denote
the trace (or \textbf{character}) of $\pi $. The following two lemmas are
well-known (due to Frobenius or Schur):

\begin{lemma}
$C(G_{*})=\mathrm{span}\{\mathrm{tr}$\medskip\ $\pi \;|\;\pi \in G^{*}\}$
and $\dim $\medskip\ $C(G_{*})=|G_{*}|=|G^{*}|$. In other words, the
characters form a basis for the vector space of class functions.
\end{lemma}

This lemma will be referred to as the \textbf{completeness lemma}.

\begin{lemma}
For $\pi ,\pi ^{\prime }\in G^{*}$, we have
\[
{1\over |G|}\sum_{g\in G}\mathrm{tr}\medskip\ \pi\, (g)\cdot \overline{\mathrm{tr}\medskip\ \pi
^{\prime }(g)}=\left\{ 
\begin{array}{cc}
1, &\pi \cong \pi ^{\prime }, \\ 
0, &\mathrm{otherwise}.
\end{array}
\right. .
\]
\end{lemma}

This lemma will be referred to as \textbf{orthogonality of characters}.

\begin{corollary} For all $\pi \in G^*$, we have
\begin{equation}
\label{eq:orthog}
{1\over |G|}\sum_{g\in G} {\rm tr}\,\pi\,(g)
=
\left\{
\begin{array}{cc}
1, & \pi\cong 1,\\
0, & {\rm otherwise},
\end{array}
\right.
\end{equation}
\end{corollary}

For $f_1,f_2\in C(G)$, let
\[
(f_1,f_2)={1\over |G|}\sum_{g\in G}f_1 (g)\cdot \overline{f_2(g)}.
\]
We call this the {\bf Schur inner product} on $C(G)$.
Note that the collection of functions
\[
\delta_x(y)=
\left\{
\begin{array}{cc}
1,& {\rm if }\ x=y,\\
0, & {\rm if}\ x\not=y.
\end{array}
\right. ,\ \ \ \ \ \ \ 
\]
for $x\in G$, 
form a basis for the vector space $C(G)$ which is orthonormal
with respect to the Schur inner product. The above lemma
says that the collection of irreducible characters form an
orthonomal basis as well.

\begin{example}
Let $G=S_3$, the symmetric group on 3 letters. Let 
$\sigma =(1\;2),\;\tau =(1\;2\;3)$, so
\[
G=\{1,\sigma ,\tau ,\tau ^2,\sigma \tau ,\sigma \tau ^2\}.
\]

Using the fact that $\sigma \tau \sigma ^{-1}=\tau ^2$, it can be shown that 
$G$ has only $3$ distinct conjugacy classes:
\[
G_{*}=\{\{1\},\{\sigma \},\{\tau \}\}.
\]

In fact, 
\[
|\{1\}|=1,\;\;\;\;\;|\{\sigma \}|=3,\;\;\;\;|\{\tau \}|=2.
\]

There are three inequivalent irreducible complex representations of $G$:
\begin{eqnarray*}
\pi _1 &:&g\longmapsto \mathrm{identity\;map\;on}\;\ccc
\mathrm{(trivial\;rep)} \\
\pi _2 &:&g\longmapsto (z\longmapsto sgn(g)\cdot z)\;\mathrm{on}\;
\ccc\mathrm{(sign\;rep)} \\
\pi _3 &:&g\longmapsto \mathrm{associated}\;3\times 3\;\mathrm{permutation\;matrix}
\ \mathrm{(standard\;rep)}\\
 & &\ \ \ \ \ \ \ \ \ \ \ \ \  {\rm acting\ on\ }
\ccc^3/\{(x,y,z)\ |\ x+y+z=0\}.
\end{eqnarray*}

Their values on the conjugacy classes are summarized as follows:
\begin{center}
\begin{tabular}{|c|c|c|c|} \hline
tr\ $\pi\backslash |\{\gamma\}|$& 1 & 3 & 2 \\ \hline
tr\ $\pi\backslash \{\gamma\}$& $\{1\}$ & $\{\sigma \}$ & $\{\tau \}$ \\ \hline
tr\ $\pi _1$ & 1 & 1 & 1 \\ 
tr\ $\pi _2$ & 1 & -1 & 1 \\ 
tr\ $\pi _3$ & 2 & 0 & -1\\ \hline
\end{tabular}
\end{center}
We shall explain why the columns of this table are orthogonal in the 
next section.

\end{example}

\begin{example}
The group $A_5$ of even permutations on $\{1,2,3,4,5\}$
has $5$ conjugacy classes:
\[
\{1\},\ \ \ \{a=(1,2)(3,4)\},\ \ \ \{b=(1,2,3)\},
\]
\[
\{c=(1,2,3,4,5)\},\ \ \ 
\{d=(1,3,5,2,4)\}.
\]
There are $5$ ireducible characters.
Their values on the conjugacy classes are summarized as follows:
\begin{center}
\begin{tabular}{|c|c|c|c|c|c|} \hline
tr\ $\pi\backslash |\{\gamma\}|$& 1 & 15 & 20 & 12 & 12 \\ \hline
tr\ $\pi\backslash \{\gamma\}$& $\{1\}$ & $\{a \}$ & $\{b \}$ & $\{ c\}$ & $\{ d\}$ \\ \hline
tr\ $\pi _1$ & 1 & 1 & 1 & 1 & 1\\ 
tr\ $\pi _2$ & 3 & -1 & 0 & $\phi$ & $\overline{\phi}$ \\ 
tr\ $\pi _3$ & 3 & -1 & 0 & $\overline{\phi}$ & $\phi$ \\
tr\ $\pi _4$ & 4 & 0 & 1 & -1 & -1 \\ 
tr\ $\pi _5$ & 5 & 1 & -1 & 0 & 0 \\ \hline
\end{tabular}
\end{center}
Here $\phi=(1+\sqrt{5})/2$ and $\overline{\phi}=(1-\sqrt{5})/2$.
\end{example}

We close this section with a

\begin{remark}
One may show (without using the completeness lemma) that the orbital 
integral map $O:C(G)\rightarrow C(G_{*})$ is surjective and that
\[
\ker \medskip\ O=\mathrm{span}\{f-f^g\;|\;f\in C(G),\;\;g\in G\},
\]
where $f^g(x)=f(g^{-1}xg)$, for $x,g\in G$.
The is equivalent to saying that the sequence

\[
0\rightarrow I(G)\stackrel{i}{\rightarrow} C(G)\stackrel{O}{\rightarrow} C(G_*)\rightarrow 0,
\]
is exact, where $I(G)=C(G)^{G-1}=\{f^g-f\ |\ f\in C(G),\ g\in G\}$
and $i$ is inclusion. 

Since these three groups $I(G)$, $C(G)$ and $C(G_*)$ are all $G$-module
(under the action induced by conjugation), one may ask what their
cohomology groups are. Since the action of $G$ on $C(G_*)$ is trivial and
since $C(G_*)$ is a finite dimensional complex vector space, we know that
$H^1(G,C(G_*))=Hom(G,C(G_*))=0$. What are the other cohomology groups?
\end{remark}

\section{A Fourier expansion}

Let 
\[
\overline{\pi}(f) = \int_G f(g)\overline{\pi(g)}\, dg
=\sum_{g\in G} f(g)\overline{\pi}(g)
\]
and
\[
{\rm tr}\,\overline{\pi}\,(f)=
\int_G {\rm tr}\,\overline{\pi}\,(g)f(g)\, dg
=\sum_{g\in G} {\rm tr}\,\overline{\pi}\,(g)f(g).
\]
By completeness, each orbital integral may be written as a linear
combination of characters. Indeed, we have the following
expansion.

\begin{lemma} For $f\in C(G)$ and $\gamma\in G$, we have
\begin{equation}
O(f,\gamma)
=\sum_{\pi\in G^*}
%\frac{1}{d_\pi^2}
 {\rm tr}\,\overline{\pi}\, (f)\cdot {\rm tr}\,\pi\,(\gamma).
\label{eq:2.1}
\end{equation}

\end{lemma}

\pf 
By the completeness lemma, there are constants
$a_\pi(f)\in \ccc$ such that
\[
O(f,\gamma)=\sum_{\pi\in G^*} a_\pi(f)\cdot {\rm tr}\,\pi\,(\gamma).
\]
Using the orthogonality of characters, we have
\[
\begin{split}
|G|\cdot a_\pi(f)
&=\int_G (\sum_{\pi'\in G^*} a_{\pi'}(f)\cdot 
{\rm tr}\,\pi'(x))\overline{{\rm tr}\,\pi\, (x)}\, dx\\
&=\int_G O(f,x)\overline{{\rm tr}\,\pi(x)}\, dx\\
&=\int_G \int_G f(g^{-1}xg)\overline{{\rm tr}\,\pi\,(x)}\, dgdx\\
&=\int_G \int_G f(x)\overline{{\rm tr}\,\pi\,(g^{-1}xg)}\, dgdx\\
&=|G|{\rm tr}\,\overline{\pi}\,(f).
\end{split}
\]
\qed

For example, let $\gamma'\in G$ and 
let $f$ denote the characteristic function of the
{\it set} $\{\gamma'\}$. Then
${\rm tr}\,{\pi}(f)=
{\rm tr}\,{\pi}(\gamma')$,
In this case,
\[
\sum_{\pi\in G^*}
 {\rm tr}\,\overline{\pi}(f)\cdot {\rm tr}\,\pi(\gamma)
=\sum_{\pi\in G^*}
 {\rm tr}\,\overline{\pi}(\gamma')\cdot {\rm tr}\,\pi(\gamma).
\]
By the lemma above, this is $=0$ if
$\gamma$ is not conjugate to $\gamma'$
(and is non-zero if $\gamma\in \{\gamma'\}$).
This is equivalent to saying that
distinct columns of the character table of a finite group are orthogonal.


\section{An analog of the global trace formula}


Let $\Gamma $ be a subgroup of $G$ and let
$G_\gamma$ denote the centralizer of $\gamma$ in $G$:
\[
G_\gamma = \{g\in G\ |\ g\gamma = \gamma g\}.
\]

Let
\[
C(\Gamma \backslash G)=\{\phi:G\rightarrow \ccc\ |\ \phi(hg)=\phi(g),
\ \forall h\in \Gamma ,\ \forall g\in G\}.
\]
Note that $G$ acts on $C(\Gamma \backslash G)$ by right translation:
\[
(R(g)\phi)(x)=\phi(xg^{-1}),\ \ \ x,g\in G.
\]
(To check this we must verify that, for each $\phi\in 
C(\Gamma \backslash G)$, we have (a) $R(1)\phi=\phi$, 
which is obvious, and (b) $R(g_1)R(g_2)\phi=R(g_1g_2)\phi$,
which follows from the equation
\[
\begin{split}
(R(g_1)R(g_2)\phi)(x)
&=(R(g_1)\phi_{g_2})(x)=\phi_{g_2}(xg_1^{-1})\\
&=\phi(xg_2^{-1}g_1^{-1})=\phi(x(g_1g_2)^{-1})=(R(g_1g_2)\phi)(x),
\ \ \ x,g_1,g_2\in G,
\end{split}
\]
where $\phi_y(x)=\phi(xy^{-1})$.) In other words, each
$g\in G$ gives rise to an automorphism
\[
R(g): C(\Gamma \backslash G)\rightarrow C(\Gamma \backslash G).
\]
This is called the {\bf right regular representation
of $G$ on $C(\Gamma \backslash G)$}.
For each fixed $f \in C(G)$, define
\[
R(f): C(\Gamma \backslash G)\rightarrow C(\Gamma \backslash G)
\]
by
\[
(R(f)\phi)(x)=\int_G f(y)(R(y)\phi)(x)\ dy,\ \ \ x\in G.
\]
This is called the {\bf right regular representation
of $C(G)$ on $C(\Gamma \backslash G)$}. We may rewrite this as
\begin{equation}
\label{eq:rt_reg}
\begin{split}
(R(f)\phi)(x)
&=\int_G f(y)(R(y)\phi)(x)\ dy\\
&=\int_G f(y)\phi(xy^{-1})\ dy\\
&=\int_G {1\over {\rm meas}(\Gamma )}\int_\Gamma  f(y)\phi(hxy^{-1})\ dh\ dy\\
&=\int_G K_f(x,y)\phi(y)\ dy,
\end{split}
\end{equation}
where
\[
K_f(x,y)={1\over {\rm meas}(\Gamma )}\int_\Gamma  f(y^{-1}hx)\ dh
\]
and where ${\rm meas}(\Gamma )=\int_\Gamma  1\ dh$ is the {\bf measure}
of $\Gamma $. The function $K_f:G\times G\rightarrow \ccc$
is called the {\bf kernel function} of the right regular representation
of $G$ on $C(\Gamma \backslash G)$. 

Choose the measure $dh$ as the counting measure,
so that ${\rm meas}(\Gamma )=|\Gamma |$. Another way to write this is
\[
K_f(x,y)={1\over |\Gamma |}\sum_{\gamma \in \Gamma } f(y^{-1}\gamma x).
\]

\begin{lemma} ${\rm tr}(R(f))={|\Gamma |\over |G|}\sum_{g\in G} K_f(x,x)$.
\end{lemma}

\pf
The delta functions $\{\delta_x\ |\ x\in \Gamma \backslash G\}$,
form an orthonormal basis for $C(\Gamma \backslash G)$ with respect to
the inner product
\[
(f_1,f_2)
={|\Gamma |\over |G|}\sum_{x\in \Gamma \backslash G}f_1 (x)\cdot 
\overline{f_2(x)}.
\]
These elements may be used to compute a matrix 
representation of $R(f)$. The trace is therefore given by
\[
\begin{split}
{\rm tr}(R(f))
&=\sum_{g\in \Gamma \backslash G}(R(f)\delta_g,\delta_g)
={|\Gamma |\over |G|}\sum_{g\in \Gamma \backslash G}
\sum_{x\in \Gamma \backslash G}(R(f)\delta_g)(x)\delta_g(x)\\
&={|\Gamma |\over |G|}\sum_{g\in \Gamma \backslash G}
\sum_{g\in \Gamma \backslash G}\sum_{y\in G}
f(y)(R(y)\delta_g)(x)\delta_g(x)\\
&={|\Gamma |\over |G|}\sum_{g\in \Gamma \backslash G}
\sum_{x\in \Gamma \backslash G}\sum_{y\in G}
f(y)\delta_g(xy^{-1})\delta_g(x)\\
&={|\Gamma |\over |G|}\sum_{g\in \Gamma \backslash G}
\sum_{x\in \Gamma \backslash G}\sum_{y\in g^{-1}\Gamma x}
f(y)\delta_g(x)\\
&={|\Gamma |\over |G|}\sum_{g\in \Gamma \backslash G}\sum_{y\in \Gamma }f(g^{-1}hg)
={1\over |G|}\sum_{g\in G}\sum_{y\in \Gamma }f(g^{-1}hg),
\end{split}
\]
as desired.   
\qed

\begin{lemma} If
\[
R=\oplus_{\pi\in G^*}m^\Gamma _\pi \pi
\]
then 

(a) ${\rm tr}(R(f))=\sum_{\pi\in G^*}m^\Gamma _\pi {\rm tr}(\pi(f))$,
for all $f\in C(G)$,

(b) $m^\Gamma _\pi \not= 0$ if and only if the restriction of
$\pi$ to $\Gamma $, $res_\Gamma ^G(\pi)$,
contains the trivial representation of $\Gamma $.
\end{lemma}

\pf
We only prove (b). Let 
$\rho\in G^*$ and $f={\rm tr}(\rho)$.
By (a) and orthogonality, ${\rm tr}(R({\rm tr}(\rho)))=m_\rho$.
By the previous lemma, we have
\[
{\rm tr}(R({\rm tr}(\rho)))
={1\over |G|}\sum_{g\in G}\sum_{y\in \Gamma }{\rm tr}(\rho)(g^{-1}hg)
=\sum_{y\in \Gamma }{\rm tr}(\rho)(h).
\]
If 
\[
Res_\Gamma ^G(\rho)=
\oplus_{\sigma\in \Gamma ^*}n_\sigma \sigma
\]
then, by orthogonality, 
\[
\sum_{h\in \Gamma }{\rm tr}(\rho)(h)=|\Gamma |n_1,
\]
where $n_1$ denotes the multiplicity of the trivial
representation of $\Gamma $ in $res_\Gamma ^G(\rho)$.
\qed

\begin{lemma}
$\int_G K_f(x,x)\, dx = {1\over |\Gamma |}
\sum_{\gamma\in \Gamma ^G} |G/G_\gamma|O(f,\gamma)$, 
where $G_\gamma$ is the centralizer of $\gamma$ and where
\[
\Gamma ^G = \{\{g^{-1}hg\ |\ g\in G\}\ |\ h\in \Gamma \}
\]
is the set of $G$-conjugacy classes in $\Gamma $.
\end{lemma}

This will be called the ``geometric side" of 
the ``global" trace formula
for finite groups. It is the trace of the
operator $R(f)$ in (\ref{eq:rt_reg}).

\pf
Since
\[
\begin{array}{c}
G/G_\gamma \rightarrow \{\gamma\} \\
g G_\gamma \longmapsto g\gamma g^{-1},
\end{array}
\]
is a bijection, we have
\[
\begin{split}
\int_G K_f(x,x)\, dx 
&= {1\over |\Gamma |}
\sum_{\gamma\in \Gamma } \int_G f(x^{-1}\gamma x)\, dx\\
&= {1\over |\Gamma |}
\sum_{\gamma\in \Gamma ^G} \sum_{\gamma'\in \{\gamma\}}
\int_G f(x^{-1}\gamma' x)\, dx\\
&= {1\over |\Gamma |}
\sum_{\gamma\in \Gamma ^G} |\{\gamma\}| O(f,\gamma)\\
&={1\over |\Gamma |}
\sum_{\gamma\in \Gamma ^G} |G/G_\gamma| O(f,\gamma).
\end{split}
\]
\qed

On the other hand, 
substituting the Fourier expansion (\ref{eq:2.1}) into the 
result of the above lemma, we obtain
\[
\begin{split}
\int_G K_f(x,x)\, dx 
&=
{1\over |\Gamma |}\sum_{\gamma\in \Gamma ^G} |G/G_\gamma| 
\sum_{\pi\in G^*} 
{\rm tr}\,\overline{\pi}(f)\cdot {\rm tr}\,\pi(\gamma)\\
&={1\over |\Gamma |}\sum_{\pi\in G^*} 
{\rm tr}\,\overline{\pi}(f)
\sum_{\gamma\in \Gamma ^G} |\{\gamma\}|{\rm tr}\,\pi(\gamma)\\
&={1\over |\Gamma |}\sum_{\pi\in G^*} 
{\rm tr}\,\overline{\pi}(f)
\sum_{g\in \Gamma } {\rm tr}\,\pi(g).
\end{split}
\]

This will be called the ``spectral side".

Setting the geometric side equal to the spectral side gives the

\begin{theorem} (``global" trace formula for finite groups)
If $R$ is the right regular representation of
$C(G)$ on $C(G/\Gamma )$ then 
\[
R=\oplus_{\pi\in G^*}m^\Gamma _\pi \pi
\]
then, for all $f\in C(G)$, we have
\[
\sum_{\pi\in G^*}m^\Gamma _\pi {\rm tr}(\pi)(f)=
\sum_{\gamma\in \Gamma ^G} |G_\gamma|^{-1}O(f,\gamma)
={1\over |G|}
\sum_{\pi\in G^*} 
{\rm tr}\,\overline{\pi}(f)\sum_{h\in \Gamma } {\rm tr}\,\pi(h).
\]
\end{theorem}

\subsection{Special cases of the ``global'' trace formula}

The following result may be regarded as an analog of
the Plancherel theorem.

\begin{corollary} Let $\Gamma =1$ and let $f\in C(G)$ be arbitrary.
Then
\[
f(1)={1\over |G|}\sum_{\pi\in G^*} 
({\rm dim}\ \pi){\rm tr}\,\overline{\pi}(f).
\]
\end{corollary}

\begin{example} Assume that $\Gamma $ is a subgroup of $G$.
If $\pi$ is an irreducible representation induced from $\Gamma $,
say $\pi=ind_\Gamma ^G\sigma$ with $\sigma\not= 1$, then
the Frobenius formula states that
\[
{\rm tr}\,\pi(g)
=\sum_{x\in G, x^{-1}gx\in \Gamma } {\rm tr}\,\sigma(x^{-1}gx).
\]
This imples 
\[
\sum_{g\in \Gamma } {\rm tr}\,\pi(g)=
\sum_{g\in \Gamma }\sum_{x\in G, x^{-1}gx\in \Gamma } {\rm tr}\,\sigma(x^{-1}gx).
\]
Orthogonality (\ref{eq:orthog}) implies
\[
\sum_{g\in \Gamma }{\rm tr}\,\sigma(x^{-1}gx)=0,
\]
since $x^{-1}\Gamma x= \Gamma $.  Therefore, such induced representations do
not contribute to the ``trace formula" above.
\end{example}

Consider the case $\Gamma =G$ of the above theorem.
Since ${\rm tr}\,\overline{1}(f)=\int_G f(g)\, dg$, the orthogonality of 
characters implies the following result.

\begin{corollary} (a) 
$\int_G f(g) dg = 
\sum_{\gamma\in G_*} |G_\gamma|^{-1} O(f,\gamma)$, for $f\in C(G)$.

(b) $\int_G f(g)\, dg= 
\sum_{\pi\in G^*}
%\frac{1}{d_\pi^2}
m_\pi^G
 {\rm tr}\,\overline{\pi}(f)$.
\end{corollary}

We call (a) the {\bf Weyl integration formula} for (the finite group) $G$.
(This formula has an easier proof than the one above: as an exercise the
reader may try to prove it using the decomposition
$G=\bigcup_{\gamma\in G_*}\{\gamma\}$). Part (b) is another form of
the Plancherel theorem.

We next consider the case where $\Gamma $ is a subgroup
of $G$ of index 2 in the above theorem. In this case, $\Gamma $ is automatically
a normal subgroup of $G$. There is a theorem of Clifford which tells
us precisely which representations in $G^*$ are induced from $\Gamma $.

\begin{lemma} (Clifford) Let $\pi\in G^*$. Fix any $s\in G-\Gamma $,
where $\Gamma $ is a subgroup of $G$ of index 2. Either

(a) the restriction of $\pi$ is irreducible, in which case $\pi$
is not induced from a representation of $\Gamma $ but
$\pi$ is an irreducible constituent of a representation 
$ind_\Gamma ^G\sigma$, for some $\sigma\in \Gamma ^*$ which satisfies 
$\sigma\cong \sigma^s$ and the restriction of $\pi$ to $\Gamma $
is $\sigma$, 

or

(b) the restriction of $\pi$ is reducible, say $\pi=\sigma \oplus \sigma'$
(some $\sigma,\sigma'\in \Gamma ^*$) in which case $\sigma'\cong\sigma^s$,
$\sigma'$ is inequivalent to $\sigma$, and
\[
\pi\cong ind_\Gamma ^G\sigma \cong ind_\Gamma ^G\sigma' .
\]
\end{lemma}

By orthogonality, the sum $\sum_{g\in \Gamma } {\rm tr}\,\pi(g)$
vanishes if the restriction of $\pi$ to $\Gamma $ is irreducible and non-trivial.
But by the remark above, the sum $\sum_{g\in \Gamma } {\rm tr}\,\pi(g)$
vanishes if $\pi$ is induced irreducibly from $\Gamma $. By Cliffords' lemma,
one of these cases must hold, so we have the following result.

\begin{lemma} Assume $\Gamma $ is a subgroup of $G$ of index $2$.
For all $f\in C(G)$, we have
\[
\sum_{\gamma\in \Gamma ^G} |G/G_\gamma|O(f,\gamma)
={1\over |\Gamma |}\sum_{\pi\in G^*, \pi|_{\Gamma }=1} {\rm tr}\,\overline{\pi}(f).
\]
\end{lemma}

\section{Analog of the local trace formula for finite groups}

We try to compute 
\[
\int_G \int_G f_1(x^{-1}yx)f_2(y)\, dxdy
\]
in two ways.

Let
\[
h(g) = \int_G f_1(x^{-1}gx)f_2(g)\, dx.
\]
The Weyl integration formula gives 
\[
\begin{split}
\int_G \int_G f_1(x^{-1}gx)f_2(g)\, dxdg
&=\int_G h(g)\, dg\\
&=\sum_{\gamma\in G_*} |G_\gamma|^{-1} O(h,\gamma)\\
&=\sum_{\gamma\in G_*} |G_\gamma|^{-1} O(f_1,\gamma)O(f_2,\gamma).
\end{split}
\]
This will be called the ``geometric side".

The Fourier expansion (\ref{eq:2.1}) gives
\[
\begin{split}
\int_G \int_G f_1(x^{-1}gx)f_2(g)\, dxdg
&=\int_G O(f_1,y)f_2(y)\, dy\\
&=\sum_{\pi\in G^*} 
{\rm tr}\,\overline{\pi}(f_1){\rm tr}\,\pi(f_2).
\end{split}
\]
This will be called the ``spectral side".

Setting the geometric side equal to the spectral side gives the

\begin{theorem} (``local" trace formula for finite groups)
For $f_1,f_2\in C(G)$, we have
\[
\sum_{\pi\in G^*} 
{\rm tr}\,\overline{\pi}(f_1){\rm tr}\,\pi(f_2)
=\sum_{\gamma\in G_*} |G_\gamma|^{-1} O(f_1,\gamma)O(f_2,\gamma).
\]
\end{theorem}

\subsection{Special cases of the local trace formula}

Setting $f_1=f_2$ gives the following curious identity:

\begin{corollary} For $f\in C(G)$ real-valued, we have
\[
\sum_{\pi\in G^*} 
%\frac{1}{d_\pi^2}
|{\rm tr}\,\pi(f)|^2
=\sum_{\gamma\in G_*} |G_\gamma|^{-1} O(f,\gamma)^2.
\]
\label{cor:squares}
\end{corollary}

\begin{example}
If $G=A_5$ and $f$ is the characteristic function
of the conjugacy class $\{(1,2,3)\}\subset A_5$
then
\footnote{This computation was performed with the aid of GAP,
a group theory software package. For more details, see
\cite{J}.}
$\sum_{\pi\in G^*} |{\rm tr}\,\pi(f)|^2
=400+400+400=1200$ and
$\sum_{\gamma\in G_*} |G_\gamma|^{-1} O(f,\gamma)^2=1200$,
as expected.
\end{example}

Next, we write the local trace formula for finite groups
down more explicitly.
First, some notation. If we let $r=|G_*|=|G^*|$ then we can let

\[
\begin{split}
G_*&=\{\gamma_1,...,\gamma_r\},\\
G^*&=\{\pi_1,...,\pi_r\},\\
G_i&=G_{\gamma_i},\ \ \ \ \ \ 1\leq i\leq r,\\
f_\gamma &= |G|^{-1}ch_\gamma,\ \ \ \ \ \gamma\in G_*,
\end{split}
\]
where $ch_X$ denotes the characteristic function of a finite set 
$X\subset G$. Finally, let
\[
a_{\ell m}={\rm tr}\,\pi_\ell(f_{\gamma_m}),\ \ \ \ 1\leq \ell,m\leq r.
\]
(These $a_{ij}$ are equal to $|G_j|^{-1}$ times the $ij^{th}$ entry
in the character table for $G$.) Then $O(\gamma_i,f_j)=\delta_{ij}$,
where $\delta_{ij}$ is the Kronecker delta function,
\[
\delta_{ij}
=\left\{
\begin{array}{cc}
1, & i=j,\\
0, &i\not= j,
\end{array}
\right.
\]
and where $a_{ij}=|G_j|^{-1}{\rm tr}\,\pi_i(\gamma_j)$.

Let $A=(a_{ij})_{1\leq i,j\leq r}=(\vec{a}_1,...,\vec{a}_r)$,
where $\vec{a}_k$ is the $k-th$ column vector of the matrix $A$.
If we put $f_1=f_{\gamma_j}$ and $f_2=f_{\gamma_i}$
then the trace formula in this case says 
\[
\langle \vec{a}_i,\vec{a}_j\rangle
=|G_i|^{-1}\delta_{ij},\ \ \ \ \ 1\leq i,j\leq r,
\]
where $\langle\vec{v},\vec{w}\rangle =\sum_i v_i\overline{w}_i$ is the usual 
Hermitian inner
product on $\ccc^r$. In other words, the trace formula in this 
case says 
\[
A^*A={\rm diag}(|G_1|^{-1},...,|G_r|^{-1}),
\]
where $A^* = \overline{A^t}$ is the conjugate transpose.

\begin{example} Let $G=S_3$ and let
\[
\gamma_1=1,\ \ \ \ \gamma_2=(1\ 2),\ \ \ \gamma_3=(1\ 2\ 3).
\]
and we label $\pi_1$, $\pi_2$, and $\pi_3$ as before. We have
\[
A=
\left(
\begin{array}{ccc}
a_{11} & a_{12} & a_{13}  \\
a_{21} & a_{22} & a_{23}  \\
a_{31} & a_{32} & a_{33}  
\end{array}
\right)
=\left(
\begin{array}{ccc}
1/6 & 1/2 & 1/3 \\
1/6 & -1/2 & 1/3 \\
1/3 & 0 & -1/3  
\end{array}
\right)
=
{1\over 6}
\left(
\begin{array}{ccc}
1 & 3 & 2 \\
1 & -3 & 2 \\
2 & 0 & -2  
\end{array}
\right)
\]
and
\[
|G_1|=6,\ \ \ \ \ |G_2|=2,\ \ \ \ |G_3|=3.
\]
It is easy to check that
\[
A\cdot A^* ={\rm diag}(6^{-1},2^{-1},3^{-1}),
\]
as predicted by the trace formula.
(However, $A$ is not normal and $A^*A$ is not
diagonal.)
\end{example}

\begin{exercise} Compute explicitly both sides of the
equation in corollary \ref{cor:squares} in case $G=S_3$ and 
$f=a_1f_{\gamma_1}+a_2f_{\gamma_2}+a_3f_{\gamma_3}$,
where $a_1,a_2,a_3$ are arbitrary real coefficients.
\end{exercise}

\begin{thebibliography}{99}  

\bibitem[GAP]{GAP} Groups, Algorithms, and Programming
(GAP - a free system for computational discrete algebra), at
\newline
\verb+ http://www.gap-system.org/+

\bibitem[J]{J} D. Joyner, ``Representations of $A_5$ using GAP,''
at
\newline
\verb+http://web.usna.navy.mil/~wdj/papers.html+

\bibitem[S]{S}
J.-P. Serre, {\bf Linear representations of finite groups}, 
Springer-Verlag, Grad Texts in Math 
\underbar{42}, 1977

\end{thebibliography}

\end{document}




